\chapter*{How to Read This Book}
\addcontentsline{toc}{chapter}{How to Read This Book}

The book is written as one argument rather than as a collection of disconnected mathematical curiosities. Chapters 1 and 2 establish the problem of descriptive language and the formal theory of persistence. Chapters 3 through 7 develop the mathematical evidence through oscillator dynamics, Platonic symmetry, the exact $\ell=5$ quartic theorem, sixth-order invariant resolution, geometry, and wave conjugation. Chapter 8 is the hinge of the book: it states the permanent separation between descriptive-language adequacy and ontological sufficiency. Chapters 9 through 11 then formulate materialist, idealist, and relational alternatives and develop the empirical consciousness program. Chapter 12 returns to the original question with the mathematical and methodological results in hand.

Readers interested primarily in the strongest exact mathematics can move directly from Chapters 4 through 7 to Appendix B. Readers interested primarily in the physical--experiential inference problem should still read Chapter 8 before Chapters 9 through 11, because the central methodological claim is that success at the descriptive level cannot be substituted for success at the ontological level. Appendix A freezes notation and the status of claims. Appendix C isolates the conditional antiunitary hypothesis so that it cannot silently become the ontology. Appendix D states defeaters and limitations in the strongest form we can currently formulate. Appendix E positions the program relative to symmetry and bifurcation theory, structural realism, physicalist consciousness theories, integrated information theory, no-report paradigms, and related philosophical traditions.

Three rules govern every chapter. Indistinguishability under a restricted descriptive algebra does not imply identity under every richer descriptive algebra. Descriptive-language adequacy does not imply ontological sufficiency. Mathematical admissibility does not imply phenomenal occupancy. These are not disclaimers added to a stronger hidden argument; they are the discipline that allows the stronger argument to remain precise.
